\clearpage
\oldpage[II]{207}
\markboth{A. GROTHENDIECK --- Chap. II}{A. GROTHENDIECK --- Chap. II}
\phantomsection
\addcontentsline{toc}{section}{Index des notations}
\section*{Index des notations}
\thispagestyle{plain}

\begingroup
\small
\setlength{\parindent}{0pt}
\setlength{\parskip}{2pt}
\newcommand{\egaShProj}{\mathcal{P}\!\operatorname{roj}}

$\mathcal{A}(X),\mathcal{A}(\mathcal{F}),\mathcal{A}(\mathcal{B})$
($X$ $S$-préschéma, $\mathcal{F}$ $\mathcal{O}_X$-Module,
$\mathcal{B}$ $\mathcal{O}_X$-Algèbre) : \hyperref[II.1.1.1-fr]{1.1.1}.

$\mathcal{A}(h)$ ($h$ $S$-morphisme) : \hyperref[II.1.1.2-fr]{1.1.2}.

$\Spec(\mathcal{B})$ ($\mathcal{B}$ $\mathcal{O}_S$-Algèbre) : \hyperref[II.1.3.1-fr]{1.3.1}.

$\widetilde{\mathcal{M}}$ ($\mathcal{M}$ $\mathcal{B}$-Module,
$\mathcal{B}$ $\mathcal{O}_S$-Algèbre) : \hyperref[II.1.4.3-fr]{1.4.3}.

$\mathbf{T}(E),\mathbf{S}(E),\mathbf{S}_A(E)$ ($E$ $A$-module) : \hyperref[II.1.7.1-fr]{1.7.1}.

$\mathbf{S}(\mathcal{E}),\mathbf{S}_{\mathcal{A}}(\mathcal{E})$
($\mathcal{E}$ $\mathcal{A}$-Module) : \hyperref[II.1.7.4-fr]{1.7.4}.

$\mathbf{V}(\mathcal{E})$ ($\mathcal{E}$ $\mathcal{O}_S$-Module quasi-cohérent) : \hyperref[II.1.7.8-fr]{1.7.8}.

$S_n,S_+,M_n,S^{(d)},M^{(d,k)},M^{(d)}$
($S$ anneau gradué, $M$ $S$-module gradué) : \hyperref[II.2.1.1-fr]{2.1.1}.

$M(n),S(n)$ ($S$ anneau gradué, $M$ $S$-module gradué) : \hyperref[II.2.1.1-fr]{2.1.1}.

$\Hom_S(M,N)$ ($S$ anneau gradué, $M,N$ $S$-modules gradués) : \hyperref[II.2.1.2-fr]{2.1.2}.

$\mathfrak N_+,\mathfrak r_+(\mathfrak J)$
($\mathfrak J$ idéal de $S_+$) : \hyperref[II.2.1.10-fr]{2.1.10}.

$S_{(f)},M_{(f)}$ ($f$ élément homogène de $S_+$) : \hyperref[II.2.2.1-fr]{2.2.1}.

$S_{(T)},M_{(T)},S_{(\mathfrak p)},M_{(\mathfrak p)}$
($T$ partie multiplicative de $S_+$ formée d'éléments homogènes,
$\mathfrak p$ idéal premier gradué de $S_+$) : \hyperref[II.2.2.7-fr]{2.2.7}.

$\Proj(S)$ ($S$ anneau gradué) : \hyperref[II.2.3.1-fr]{2.3.1}.

$V_+(E)$ ($E$ partie de $S_+$) : \hyperref[II.2.3.2-fr]{2.3.2}.

$D_+(f)$ ($f$ élément de $S$) : \hyperref[II.2.3.3-fr]{2.3.3}.

$\mathfrak j_+(Y)$ ($Y$ partie de $\Proj(S)$) : \hyperref[II.2.3.10-fr]{2.3.10}.

$\widetilde M$ ($M$ $S$-module gradué) : \hyperref[II.2.5.2-fr]{2.5.2}.

$\widetilde u$ ($u$ homomorphisme de degré zéro) : \hyperref[II.2.5.4-fr]{2.5.4}.

$\mathcal{O}_X(n),\mathcal{F}(n)$
($X=\Proj(S)$, $\mathcal{F}$ $\mathcal{O}_X$-Module) : \hyperref[II.2.5.10-fr]{2.5.10}.

$\lambda$ (homomorphisme canonique) : \hyperref[II.2.5.11-fr]{2.5.11}.

$\mu$ (homomorphisme canonique) : \hyperref[II.2.5.12-fr]{2.5.12}.

$X^{(d)}$ ($X=\Proj(S)$, $d$ entier $>0$) : \hyperref[II.2.5.16-fr]{2.5.16}.

$\Gamma_*(\mathcal{F})$ ($\mathcal{F}$ $\mathcal{O}_X$-Module,
$X=\Proj(S)$) : \hyperref[II.2.6.1-fr]{2.6.1}.

$\alpha_n,\alpha,\alpha_{\mathcal{M}}$ (homomorphismes canoniques) : \hyperref[II.2.6.2-fr]{2.6.2}.

$\beta,\beta_{\mathcal{F}}$ (homomorphismes canoniques) : \hyperref[II.2.6.4-fr]{2.6.4}.

$G(\vphi),{}^a\vphi$ (par abus de langage)
($\vphi$ homomorphisme d'anneaux gradués) : \hyperref[II.2.8.1-fr]{2.8.1}.

$\Proj(\vphi)$ ($\vphi$ homomorphisme d'anneaux gradués) : \hyperref[II.2.8.2-fr]{2.8.2}.

$\mathcal{S}^{(d)},\mathcal{M}^{(d,k)},\mathcal{M}^{(d)},\mathcal{M}(n)$
($\mathcal{S}$ $\mathcal{O}_Y$-Algèbre graduée,
$\mathcal{M}$ $\mathcal{S}$-Module gradué) : \hyperref[II.3.1.1-fr]{3.1.1}.

$\Proj(\mathcal{S})$ ($\mathcal{S}$ $\mathcal{O}_Y$-Algèbre graduée quasi-cohérente) : \hyperref[II.3.1.3-fr]{3.1.3}.

$X_f$ ($X=\Proj(\mathcal{S})$, $f$ section de $\mathcal{S}_d$ au-dessus de $Y$) : \hyperref[II.3.1.4-fr]{3.1.4}.

$\widetilde{\mathcal{M}}$ ($\mathcal{M}$ $\mathcal{S}$-Module gradué) : \hyperref[II.3.2.2-fr]{3.2.2}.

$\mathcal{O}_X(n),\mathcal{F}(n)$
($X=\Proj(\mathcal{S})$, $\mathcal{F}$ $\mathcal{O}_X$-Module) : \hyperref[II.3.2.5-fr]{3.2.5}.

$\lambda,\mu$ (homomorphismes canoniques) : \hyperref[II.3.2.6-fr]{3.2.6}.

$\Gamma_*(\mathcal{F})$ ($\mathcal{F}$ $\mathcal{O}_X$-Module,
$X=\Proj(\mathcal{S})$) : \hyperref[II.3.3.1-fr]{3.3.1}.

$\alpha_n,\alpha,\alpha_{\mathcal{M}}$ (homomorphismes canoniques) : \hyperref[II.3.3.2-fr]{3.3.2}.

\oldpage[II]{208}
\markboth{A. GROTHENDIECK --- Chap. II}{A. GROTHENDIECK --- Chap. II}
$\beta,\beta_{\mathcal{F}}$ (homomorphismes canoniques) : \hyperref[II.3.3.4-fr]{3.3.4}.

$G(\vphi),\Proj(\vphi)$
($\vphi$ homomorphisme de $\mathcal{O}_Y$-Algèbres graduées) : \hyperref[II.3.5.1-fr]{3.5.1}.

$\Proj(u)$ ($u$ $g$-morphisme d'une $\mathcal{O}_Y$-Algèbre graduée
dans une $\mathcal{O}_{Y'}$-Algèbre graduée) : \hyperref[II.3.5.6-fr]{3.5.6}.

$r_{\mathcal{L},\psi}$ ($\mathcal{L}$ $\mathcal{O}_X$-Module inversible) : \hyperref[II.3.7.1-fr]{3.7.1}.

$\mathbf{P}(\mathcal{E}),\mathbf{P}(E),\mathbf{P}_Y^n,\mathbf{P}_A^n$ : \hyperref[II.4.1.1-fr]{4.1.1}.

$\mathbf{P}(u)$ ($u$ homomorphisme surjectif de $\mathcal{O}_Y$-Modules) : \hyperref[II.4.1.2-fr]{4.1.2}.

$\Hyp_Y(X,\mathcal{E})$ ($\mathcal{E}$ $\mathcal{O}_Y$-Module quasi-cohérent) : \hyperref[II.4.2.5-fr]{4.2.5}.

$\varsigma$ (morphisme de Segre) : \hyperref[II.4.3.1-fr]{4.3.1}.

$\mathcal{F}(n)$ ($\mathcal{F}$ $\mathcal{O}_X$-Module, $X$ muni d'un
$\mathcal{O}_X$-Module inversible) : \hyperref[II.4.5.1-fr]{4.5.1}.

$\varepsilon$ (automorphisme identique) : \hyperref[II.4.5.1-fr]{4.5.1}.

$P(u,T),\sigma_i(u)$ ($u$ endomorphisme de module libre de base finie) : \hyperref[II.6.4.1-fr]{6.4.1}.

$P(u,T),\sigma_i(u),\det u$
($u$ endomorphisme de module de type fini sur un anneau intègre ou un
anneau noethérien réduit) : \hyperref[II.6.4.2-fr]{6.4.2} et \hyperref[II.6.4.7-fr]{6.4.7}.

$\sigma_i(u),\det u$ ($u$ endomorphisme de $\mathcal{A}$-Module localement libre) : \hyperref[II.6.4.8-fr]{6.4.8}.

$\sigma_i,\det$ (homomorphismes de faisceaux) : \hyperref[II.6.4.8-fr]{6.4.8}, \hyperref[II.6.4.9-fr]{6.4.9} et \hyperref[II.6.4.10-fr]{6.4.10}.

$\sigma_i(u),\det u$
($u$ endomorphisme de $\mathcal{O}_X$-Module de type fini sur un préschéma
localement intègre ou un préschéma localement noethérien réduit) :
\hyperref[II.6.4.9-fr]{6.4.9} et \hyperref[II.6.4.10-fr]{6.4.10}.

$N_{\mathcal{B}/\mathcal{A}}(f)$
($f$ section de la $\mathcal{B}$-Algèbre $\mathcal{A}$) : \hyperref[II.6.5.1-fr]{6.5.1}.

$N_{\mathcal{B}/\mathcal{A}}(\mathcal{L}')$
($\mathcal{L}'$ $\mathcal{B}$-Module inversible) : \hyperref[II.6.5.2-fr]{6.5.2}.

$N_{\mathcal{B}/\mathcal{A}}(h')$
($h'$ homomorphisme de $\mathcal{B}$-modules inversibles) : \hyperref[II.6.5.3-fr]{6.5.3}.

$N_{X'/X}(\mathcal{L}')$
($X'$ fini au-dessus de $X$, $\mathcal{L}'$ $\mathcal{O}_{X'}$-Module inversible) : \hyperref[II.6.5.5-fr]{6.5.5}.

$N_{X'/X}(h')$
($h'$ homomorphisme de $\mathcal{O}_{X'}$-Modules inversibles) : \hyperref[II.6.5.5-fr]{6.5.5}.

$\mathcal{S}_X$ ($X=\Proj(\mathcal{S})$, $\mathcal{S}$ somme d'Idéaux de
$\mathcal{O}_Y$) : \hyperref[II.8.1.5-fr]{8.1.5}.

$S^\geq,S^\leq,S_f^\geq,S_f^\leq,
M^\geq,M^\leq,M_f^\geq,M_f^\leq$
($S$ anneau gradué, $M$ $S$-module gradué,
$f$ élément homogène de $S_+$) : \hyperref[II.8.2.1-fr]{8.2.1}.

$\widehat S$ ($S$ anneau gradué) : \hyperref[II.8.2.2-fr]{8.2.2}.

$\widehat M$ ($M$ $S$-module gradué) : \hyperref[II.8.2.4-fr]{8.2.4}.

$S_{[n]},S^\natural,f^\natural$
($S$ anneau gradué, $f$ homogène dans $S$) : \hyperref[II.8.2.6-fr]{8.2.6}.

$M_{[n]},M^\natural$ ($M$ $S$-module gradué) : \hyperref[II.8.2.8-fr]{8.2.8}.

$\widehat{\mathcal{S}}$ ($\mathcal{S}$ $\mathcal{O}_Y$-Algèbre graduée) : \hyperref[II.8.3.1-fr]{3.8.1}.

$G_m$ (schéma en groupes) : \hyperref[II.8.3.9-fr]{8.3.9}.

$\mathcal{S}_X$ ($X=\Proj(\mathcal{S})$, $\mathcal{S}$ $\mathcal{O}_Y$-Algèbre graduée) : \hyperref[II.8.6.1-fr]{8.6.1}.

$\mathcal{S}_X,C_X,\widehat C_X,E_X,\widehat E_X$
($X=\Proj(\mathcal{S})$, $\mathcal{S}$ $\mathcal{O}_Y$-Algèbre graduée
avec $\mathcal{S}_0=\mathcal{O}_Y$) : \hyperref[II.8.6.1-fr]{8.6.1}.

$\mathcal{S}_{[n]},\mathcal{S}^{\natural}$
($\mathcal{S}$ $\mathcal{O}_Y$-Algèbre graduée) : \hyperref[II.8.7.2-fr]{8.7.2}.

$\egaShProj_0(\mathcal{M}),\egaShProj(\mathcal{M}),\mathcal{M}_X,
\widehat{\mathcal{M}},\mathcal{M}^{\square}$
($\mathcal{M}$ $\mathcal{S}$-Module gradué) : \hyperref[II.8.12.1-fr]{8.12.1}.

$\mathcal{M}_X^\geq$
($X=\Proj(\mathcal{S})$, $\mathcal{M}$ $\mathcal{S}$-Module gradué) : \hyperref[II.8.12.5-fr]{8.12.5}.

$\mathcal{M}_{[n]},\mathcal{M}^{\natural}$
($\mathcal{M}$ $\mathcal{S}$-Module gradué) : \hyperref[II.8.12.9-fr]{8.12.9}.

$(\widetilde{\mathcal{N}})^-$
($\mathcal{N}$ sous-$\mathcal{S}$-Module d'un $\mathcal{S}$-Module gradué) : \hyperref[II.8.13.1-fr]{8.13.1}.

$\overline{\mathcal{N}}$
($\mathcal{N}$ sous-$\mathcal{S}$-Module d'un $\mathcal{S}$-Module gradué) : \hyperref[II.8.13.2-fr]{8.13.2}.

$\Gamma_*(\mathcal{F})$
($X=\Proj(\mathcal{S})$, $\mathcal{F}$ $\mathcal{S}_X$-Module gradué) : \hyperref[II.8.14.5-fr]{8.14.5}.

$\mathcal{S}_{X'},\Gamma'_*(\mathcal{F}'),\egaShProj'(\mathcal{M}),
\alpha',\beta'$
($X'$ ouvert de $X=\Proj(\mathcal{S})$, $\mathcal{F}'$ $\mathcal{S}_{X'}$-Module gradué,
$\mathcal{M}$ $\mathcal{S}$-Module gradué) : \hyperref[II.8.14.6-fr]{8.14.6}.

\endgroup

\clearpage
\oldpage[II]{209}
\markboth{A. GROTHENDIECK --- Chap. II}{A. GROTHENDIECK --- Chap. II}
\thispagestyle{plain}
\phantomsection
\addcontentsline{toc}{section}{Index terminologique}
\section*{Index terminologique}

\begingroup
\small
\setlength{\parindent}{0pt}
\setlength{\parskip}{2pt}
\newcommand{\egaTerme}[2]{\textbf{#1} : #2.\par}

\egaTerme{Affine (morphisme)}{\hyperref[II.1.6.1-fr]{1.6.1}}
\egaTerme{Affine (préschéma) au-dessus d'un préschéma}{\hyperref[II.1.2.1-fr]{1.2.1}}
\egaTerme{Algèbre graduée sur un anneau gradué}{\hyperref[II.2.1.2-fr]{2.1.2}}
\egaTerme{Algèbre symétrique d'un $A$-module}{\hyperref[II.1.7.1-fr]{1.7.1}}
\egaTerme{$\mathcal{A}$-Algèbre symétrique d'un $\mathcal{A}$-Module}{\hyperref[II.1.7.4-fr]{1.7.4}}
\egaTerme{$\mathcal{O}_Y$-Algèbre graduée essentiellement réduite}{\hyperref[II.3.1.12-fr]{3.1.12}}
\egaTerme{$\mathcal{O}_Y$-Algèbre graduée intègre}{\hyperref[II.3.1.12-fr]{3.1.12}}
\egaTerme{Ample ($\mathcal{O}_X$-Module inversible)}{\hyperref[II.4.5.3-fr]{4.5.3}}
\egaTerme{Ample relativement à $f$, $f$-ample, $Y$-ample ($\mathcal{O}_X$-Module inversible)}{\hyperref[II.4.6.1-fr]{4.6.1}}
\egaTerme{Anneau gradué essentiellement intègre}{\hyperref[II.2.1.11-fr]{2.1.11}}
\egaTerme{Anneau gradué essentiellement réduit}{\hyperref[II.2.1.10-fr]{2.1.10}}
\egaTerme{Associé à un $\mathcal{B}$-Module (faisceau)}{\hyperref[II.1.4.3-fr]{1.4.3}}
\egaTerme{Associé à un $S$-module gradué (faisceau)}{\hyperref[II.2.5.3-fr]{2.5.3}}
\egaTerme{Associé à un $\mathcal{S}$-Module gradué (faisceau)}{\hyperref[II.3.2.2-fr]{3.2.2}}
\egaTerme{Associé à une $\mathcal{O}_S$-Algèbre ($S$-schéma)}{\hyperref[II.1.3.1-fr]{1.3.1}}

\egaTerme{Condition (TF), condition (TN)}{\hyperref[II.2.7.2-fr]{2.7.2}}
\egaTerme{Condition (TF), condition (TN)}{\hyperref[II.3.4.2-fr]{3.4.2}}
\egaTerme{Cône affine, cône projectif défini par une $\mathcal{O}_Y$-Algèbre graduée}{\hyperref[II.8.3.1-fr]{8.3.1}}
\egaTerme{Cône affine épointé, cône projectif épointé défini par une $\mathcal{O}_Y$-Algèbre graduée}{\hyperref[II.8.3.4-fr]{8.3.4}}
\egaTerme{Cône projetant affine, cône projetant projectif d'un préschéma $\Proj(\mathcal{S})$}{\hyperref[II.8.3.1-fr]{8.3.1}}
\egaTerme{Courbe algébrique sur un corps}{\hyperref[II.7.4.2-fr]{7.4.2}}
\egaTerme{Courbe algébrique complète}{\hyperref[II.7.4.7-fr]{7.4.7}}
\egaTerme{Critère de Serre}{\hyperref[II.5.2.1-fr]{5.2.1}}

\egaTerme{Déterminant d'un endomorphisme d'un module de type fini sur un anneau intègre}{\hyperref[II.6.4.2-fr]{6.4.2}}
\egaTerme{Déterminant d'un endomorphisme d'un module de type fini sur un anneau noethérien réduit}{\hyperref[II.6.4.7-fr]{6.4.7}}
\egaTerme{Déterminant d'un endomorphisme d'un $\mathcal{A}$-Module localement libre}{\hyperref[II.6.4.8-fr]{6.4.8}}
\egaTerme{Déterminant d'un endomorphisme d'un $\mathcal{O}_X$-Module de type fini sur un préschéma $X$ localement intègre}{\hyperref[II.6.4.9-fr]{6.4.9}}
\egaTerme{Déterminant d'un endomorphisme d'un $\mathcal{O}_X$-Module de type fini sur un préschéma $X$ localement noethérien et réduit}{\hyperref[II.6.4.10-fr]{6.4.10}}

\egaTerme{Éclaté (préschéma)}{\hyperref[II.8.1.3-fr]{8.1.3}}
\egaTerme{Entier (morphisme)}{\hyperref[II.6.1.1-fr]{6.1.1}}
\egaTerme{Entier (préschéma) au-dessus d'un préschéma}{\hyperref[II.6.1.1-fr]{6.1.1}}
\egaTerme{Entière ($\mathcal{O}_S$-Algèbre quasi-cohérente)}{\hyperref[II.6.1.2-fr]{6.1.2}}
\egaTerme{Entière (section d'une $\mathcal{A}$-Algèbre) sur $\mathcal{A}$}{\hyperref[II.6.3.1-fr]{6.3.1}}

\egaTerme{Faisceau fondamental d'un fibré projectif}{\hyperref[II.4.1.1-fr]{4.1.1}}
\egaTerme{Fermeture intégrale d'un faisceau d'anneaux $\mathcal{A}$ dans une $\mathcal{A}$-Algèbre}{\hyperref[II.6.3.2-fr]{6.3.2}}
\egaTerme{Fermeture intégrale d'un préschéma $X$ relativement à une $\mathcal{O}_X$-Algèbre}{\hyperref[II.6.3.4-fr]{6.3.4}}
\egaTerme{Fermeture projective d'un cône affine}{\hyperref[II.8.3.1-fr]{8.3.1}}

\oldpage[II]{210}
\markboth{A. GROTHENDIECK --- Chap. II}{A. GROTHENDIECK --- Chap. II}
\egaTerme{Fibre géométrique rationnelle d'un fibré projectif sur une extension $K$ de $k(y)$}{\hyperref[II.4.2.6-fr]{4.2.6}}
\egaTerme{Fibré projectif défini par un $\mathcal{O}_Y$-Module}{\hyperref[II.4.1.1-fr]{4.1.1}}
\egaTerme{Fibré vectoriel défini par un $\mathcal{O}_S$-Module}{\hyperref[II.1.7.8-fr]{1.7.8}}
\egaTerme{Fini (morphisme)}{\hyperref[II.6.1.1-fr]{6.1.1}}
\egaTerme{Fini (préschéma) au-dessus d'un préschéma}{\hyperref[II.6.1.1-fr]{6.1.1}}
\egaTerme{Finie ($\mathcal{O}_S$-Algèbre quasi-cohérente)}{\hyperref[II.6.1.2-fr]{6.1.2}}
\egaTerme{Fonctions symétriques élémentaires d'un endomorphisme d'un module localement libre}{\hyperref[II.6.4.1-fr]{6.4.1}}
\egaTerme{Fonctions symétriques élémentaires d'un endomorphisme d'un module de type fini sur un anneau intègre}{\hyperref[II.6.4.2-fr]{6.4.2}}
\egaTerme{Fonctions symétriques élémentaires d'un endomorphisme d'un module de type fini sur un anneau noethérien réduit}{\hyperref[II.6.4.7-fr]{6.4.7}}
\egaTerme{Fonctions symétriques élémentaires d'un endomorphisme d'un $\mathcal{A}$-Module localement libre de rang constant}{\hyperref[II.6.4.8-fr]{6.4.8}}
\egaTerme{Fonctions symétriques élémentaires d'un endomorphisme d'un $\mathcal{O}_X$-Module de type fini sur un préschéma $X$ localement intègre}{\hyperref[II.6.4.9-fr]{6.4.9}}
\egaTerme{Fonctions symétriques élémentaires d'un endomorphisme d'un $\mathcal{O}_X$-Module de type fini sur un préschéma $X$ localement noethérien réduit}{\hyperref[II.6.4.10-fr]{6.4.10}}
\egaTerme{Fonctions symétriques élémentaires d'une section d'une $\mathcal{A}$-Algèbre}{\hyperref[II.6.5.1-fr]{6.5.1}}

\egaTerme{Homogénisé d'un sous-$\mathcal{S}$-Module d'un $\mathcal{S}$-Module gradué}{\hyperref[II.8.13.2-fr]{8.13.2}}
\egaTerme{Homomorphisme d'anneaux gradués}{\hyperref[II.2.1.2-fr]{2.1.2}}
\egaTerme{Homomorphisme de degré $k$ de modules gradués}{\hyperref[II.2.1.2-fr]{2.1.2}}

\egaTerme{Idéal de $S_+$, idéal premier gradué de $S_+$ ($S$ anneau gradué)}{\hyperref[II.2.1.10-fr]{2.1.10}}

\egaTerme{Lemme de Chow}{\hyperref[II.5.6.1-fr]{5.6.1}}
\egaTerme{Lieu à l'infini d'un cône projectif}{\hyperref[II.8.3.3-fr]{8.3.3}}

\egaTerme{Morphisme canonique $G(\mathcal{E})\to\Proj(S)$, $S=\bigoplus_{n\geq0}\Gamma(X,\mathcal{L}^{\otimes n})$}{\hyperref[II.4.5.1-fr]{4.5.1}}
\egaTerme{Morphisme canonique $X\to\Spec(\Gamma(X,\mathcal{O}_X))$}{\hyperref[II.5.1.1-fr]{5.1.1}}
\egaTerme{Morphisme de Segre}{\hyperref[II.4.3.1-fr]{4.3.1}}

\egaTerme{Nilradical de $S_+$ ($S$ anneau gradué)}{\hyperref[II.2.1.10-fr]{2.1.10}}
\egaTerme{Nilradical de $\mathcal{S}_+$ ($\mathcal{S}$ $\mathcal{O}_Y$-Algèbre graduée quasi-cohérente)}{\hyperref[II.3.1.12-fr]{3.1.12}}
\egaTerme{Norme d'une section d'une $\mathcal{A}$-Algèbre}{\hyperref[II.6.5.1-fr]{6.5.1}}
\egaTerme{Norme d'un $\mathcal{B}$-Module inversible}{\hyperref[II.6.5.2-fr]{6.5.2}}
\egaTerme{Norme d'une section d'un $\mathcal{B}$-Module inversible}{\hyperref[II.6.5.3-fr]{6.5.3}}
\egaTerme{Norme d'un $\mathcal{O}_{X'}$-Module inversible}{\hyperref[II.6.5.5-fr]{6.5.5}}

\egaTerme{Périodique (Algèbre graduée)}{\hyperref[II.8.14.12-fr]{8.14.12}}
\egaTerme{Polynôme caractéristique d'un endomorphisme d'un module de type fini sur un anneau intègre}{\hyperref[II.6.4.2-fr]{6.4.2}}
\egaTerme{Polynôme caractéristique d'un endomorphisme d'un module de type fini sur un anneau noethérien réduit}{\hyperref[II.6.4.7-fr]{6.4.7}}
\egaTerme{Préschéma des sommets d'un cône affine (projectif)}{\hyperref[II.8.3.3-fr]{8.3.3}}
\egaTerme{Présentation finie (module gradué ayant une)}{\hyperref[II.2.1.1-fr]{2.1.1}}
\egaTerme{Présentation finie ($\mathcal{S}$-Module gradué ayant une)}{\hyperref[II.3.1.1-fr]{3.1.1}}
\egaTerme{Produit tensoriel gradué de deux modules gradués}{\hyperref[II.2.1.2-fr]{2.1.2}}
\egaTerme{Projectif (morphisme)}{\hyperref[II.5.5.2-fr]{5.5.2}}
\egaTerme{Projectif (préschéma) au-dessus d'un préschéma}{\hyperref[II.5.5.2-fr]{5.5.2}}
\egaTerme{Propre (morphisme)}{\hyperref[II.5.4.1-fr]{5.4.1}}
\egaTerme{Propre (partie)}{\hyperref[II.5.4.10-fr]{5.4.10}}
\egaTerme{Propre (préschéma) au-dessus d'un préschéma}{\hyperref[II.5.4.1-fr]{5.4.1}}

\egaTerme{Quasi-affine (morphisme)}{\hyperref[II.5.1.1-fr]{5.1.1}}
\egaTerme{Quasi-affine (préschéma) au-dessus d'un préschéma}{\hyperref[II.5.1.1-fr]{5.1.1}}
\egaTerme{Quasi-affine (schéma)}{\hyperref[II.5.1.1-fr]{5.1.1}}

\oldpage[II]{211}
\markboth{CHAPITRE II}{CHAPITRE II}
\egaTerme{Quasi-fini (morphisme)}{\hyperref[II.6.2.3-fr]{6.2.3}}
\egaTerme{Quasi-fini (préschéma) au-dessus d'un préschéma}{\hyperref[II.6.2.3-fr]{6.2.3}}
\egaTerme{Quasi-projectif (morphisme)}{\hyperref[II.5.3.1-fr]{5.3.1}}
\egaTerme{Quasi-projectif (préschéma) au-dessus d'un préschéma}{\hyperref[II.5.3.1-fr]{5.3.1}}

\egaTerme{Racine d'un idéal de $S_+$ ($S$ anneau gradué)}{\hyperref[II.2.1.10-fr]{2.1.10}}
\egaTerme{Rétraction canonique d'un cône projectif épointé sur son lieu à l'infini}{\hyperref[II.8.3.5-fr]{8.3.5}}

\egaTerme{Section canonique de $\mathcal{O}_X(-1)$ ($X$ préschéma éclaté)}{\hyperref[II.8.1.9-fr]{8.1.9}}
\egaTerme{Section nulle d'un cône affine}{\hyperref[II.8.3.3-fr]{8.3.3}}
\egaTerme{Section nulle d'un fibré vectoriel}{\hyperref[II.1.7.9-fr]{1.7.9}}
\egaTerme{Section sommet d'un cône affine (projectif)}{\hyperref[II.8.3.3-fr]{8.3.3}}
\egaTerme{Spectre d'une $\mathcal{O}_S$-Algèbre}{\hyperref[II.1.3.1-fr]{1.3.1}}
\egaTerme{Spectre premier homogène d'un anneau gradué}{\hyperref[II.2.3.1-fr]{2.3.1}}
\egaTerme{Spectre homogène d'une $\mathcal{O}_S$-Algèbre graduée quasi-cohérente}{\hyperref[II.3.1.3-fr]{3.1.3}}

\egaTerme{(TN)-injectif, (TN)-surjectif, (TN)-bijectif (homomorphisme de modules gradués)}{\hyperref[II.2.7.2-fr]{2.7.2}}
\egaTerme{(TN)-isomorphisme de modules gradués}{\hyperref[II.2.7.2-fr]{2.7.2}}
\egaTerme{(TN)-injectif, (TN)-surjectif, (TN)-bijectif (homomorphisme d'anneaux gradués)}{\hyperref[II.2.9.1-fr]{2.9.1}}
\egaTerme{(TN)-isomorphisme d'anneaux gradués}{\hyperref[II.2.9.1-fr]{2.9.1}}
\egaTerme{(TN)-injectif, (TN)-surjectif, (TN)-bijectif (homomorphisme de $\mathcal{S}$-Modules gradués)}{\hyperref[II.3.4.2-fr]{3.4.2}}
\egaTerme{(TN)-isomorphisme de $\mathcal{S}$-Modules gradués}{\hyperref[II.3.4.2-fr]{3.4.2}}
\egaTerme{(TN)-injectif, (TN)-surjectif, (TN)-bijectif (homomorphisme de $\mathcal{O}_Y$-Algèbres graduées)}{\hyperref[II.3.6.1-fr]{3.6.1}}
\egaTerme{(TN)-isomorphisme de $\mathcal{O}_Y$-Algèbres graduées}{\hyperref[II.3.6.1-fr]{3.6.1}}
\egaTerme{Topologie spectrale sur $\Proj(S)$}{\hyperref[II.2.3.3-fr]{2.3.3}}
\egaTerme{Très ample pour $q$, très ample pour $Y$ ($\mathcal{O}_X$-Module inversible)}{\hyperref[II.4.4.2-fr]{4.4.2}}
\egaTerme{Type fini ($\mathcal{S}$-Module gradué de)}{\hyperref[II.3.1.1-fr]{3.1.1}}

\egaTerme{Universellement fermé (morphisme)}{\hyperref[II.5.4.9-fr]{5.4.9}}

\endgroup

\clearpage
\oldpage[II]{212}
\thispagestyle{empty}
\mbox{}

\clearpage
\oldpage[II]{213}
\markboth{TABLE DES MATIÈRES}{TABLE DES MATIÈRES}
\phantomsection
\addcontentsline{toc}{section}{Table des matières originale}
\section*{Table des matières}
\thispagestyle{plain}

\begingroup
\small
\renewcommand{\arraystretch}{1.04}
\begin{longtable}{@{}p{0.09\textwidth}p{0.69\textwidth}r@{}}
\textbf{Unité} & \textbf{Titre} & \textbf{Page originale}\\
\hline
\textbf{II} & \textbf{Étude globale élémentaire de quelques classes de morphismes} & \textbf{5}\\
\hyperref[section:II.1-fr]{\textsection1} & Morphismes affines & 5\\
\hyperref[subsection:II.1.1-fr]{1.1} & $S$-préschémas et $\mathcal{O}_S$-Algèbres & 5\\
\hyperref[subsection:II.1.2-fr]{1.2} & Préschémas affines sur un préschéma & 6\\
\hyperref[subsection:II.1.3-fr]{1.3} & Préschéma affine au-dessus de $S$ associé à une $\mathcal{O}_S$-Algèbre & 8\\
\hyperref[subsection:II.1.4-fr]{1.4} & Faisceaux quasi-cohérents sur un préschéma affine au-dessus de $S$ & 9\\
\hyperref[II.1.5-fr]{1.5} & Changement du préschéma de base & 12\\
\hyperref[II.1.6-fr]{1.6} & Morphismes affines & 14\\
\hyperref[II.1.7-fr]{1.7} & Fibré vectoriel associé à un faisceau de modules & 14\\
\hyperref[section:II.2-fr]{\textsection2} & Spectres premiers homogènes & 19\\
\hyperref[subsection:II.2.1-fr]{2.1} & Généralités sur les anneaux et modules gradués & 19\\
\hyperref[subsection:II.2.2-fr]{2.2} & Anneaux de fractions d'un anneau gradué & 23\\
\hyperref[subsection:II.2.3-fr]{2.3} & Spectre premier homogène d'un anneau gradué & 25\\
\hyperref[subsection:II.2.4-fr]{2.4} & La structure de schéma sur $\Proj(S)$ & 28\\
\hyperref[subsection:II.2.5-fr]{2.5} & Faisceau associé à un module gradué & 30\\
\hyperref[subsection:II.2.6-fr]{2.6} & $S$-module gradué associé à un faisceau sur $\Proj(S)$ & 36\\
\hyperref[subsection:II.2.7-fr]{2.7} & Conditions de finitude & 38\\
\hyperref[subsection:II.2.8-fr]{2.8} & Comportements fonctoriels & 41\\
\hyperref[subsection:II.2.9-fr]{2.9} & Sous-préschémas fermés d'un schéma $\Proj(S)$ & 48\\
\hyperref[section:II.3-fr]{\textsection3} & Spectre homogène d'un faisceau d'algèbres graduées & 49\\
\hyperref[subsection:II.3.1-fr]{3.1} & Spectre homogène d'une $\mathcal{O}_Y$-Algèbre graduée quasi-cohérente & 49\\
\hyperref[subsection:II.3.2-fr]{3.2} & Faisceau sur $\Proj(\mathcal{S})$ associé à un $\mathcal{S}$-Module gradué & 54\\
\hyperref[subsection:II.3.3-fr]{3.3} & $\mathcal{S}$-Module gradué associé à un faisceau sur $\Proj(\mathcal{S})$ & 56\\
\hyperref[subsection:II.3.4-fr]{3.4} & Conditions de finitude & 59\\
\hyperref[subsection:II.3.5-fr]{3.5} & Comportements fonctoriels & 61\\
\hyperref[subsection:II.3.6-fr]{3.6} & Sous-préschémas fermés d'un préschéma $\Proj(\mathcal{S})$ & 64\\
\hyperref[subsection:II.3.7-fr]{3.7} & Morphismes d'un préschéma dans un spectre homogène & 65\\
\hyperref[subsection:II.3.8-fr]{3.8} & Critères d'immersion dans un spectre homogène & 69\\
\end{longtable}
\oldpage[II]{214}
\markboth{A. GROTHENDIECK --- Chap. II}{A. GROTHENDIECK --- Chap. II}
\begin{longtable}{@{}p{0.09\textwidth}p{0.69\textwidth}r@{}}
\hyperref[section:II.4-fr]{\textsection4} & Fibrés projectifs. Faisceaux amples & 71\\
\hyperref[subsection:II.4.1-fr]{4.1} & Définition des fibrés projectifs & 71\\
\hyperref[subsection:II.4.2-fr]{4.2} & Morphismes d'un préschéma dans un fibré projectif & 72\\
\hyperref[subsection:II.4.3-fr]{4.3} & Le morphisme de Segre & 76\\
\hyperref[subsection:II.4.4-fr]{4.4} & Immersions dans les fibrés projectifs. Faisceaux très amples & 78\\
\hyperref[subsection:II.4.5-fr]{4.5} & Faisceaux amples & 83\\
\hyperref[II.4.6-fr]{4.6} & Faisceaux relativement amples & 89\\
\hyperref[section:II.5-fr]{\textsection5} & Morphismes quasi-affines; morphismes quasi-projectifs; morphismes propres; morphismes projectifs & 94\\
\hyperref[subsection:II.5.1-fr]{5.1} & Morphismes quasi-affines & 94\\
\hyperref[subsection:II.5.2-fr]{5.2} & Le critère de Serre & 97\\
\hyperref[subsection:II.5.3-fr]{5.3} & Morphismes quasi-projectifs & 99\\
\hyperref[subsection:II.5.4-fr]{5.4} & Morphismes propres et morphismes universellement fermés & 100\\
\hyperref[subsection:II.5.5-fr]{5.5} & Morphismes projectifs & 103\\
\hyperref[subsection:II.5.6-fr]{5.6} & Le lemme de Chow & 106\\
\hyperref[section:II.6-fr]{\textsection6} & Morphismes entiers et morphismes finis & 110\\
\hyperref[subsection:II.6.1-fr]{6.1} & Préschémas entiers sur un autre & 110\\
\hyperref[subsection:II.6.2-fr]{6.2} & Morphismes quasi-finis & 114\\
\hyperref[subsection:II.6.3-fr]{6.3} & Fermeture intégrale d'un préschéma & 116\\
\hyperref[subsection:II.6.4-fr]{6.4} & Déterminant d'un endomorphisme de $\mathcal{O}_X$-Module & 120\\
\hyperref[subsection:II.6.5-fr]{6.5} & Norme d'un faisceau inversible & 125\\
\hyperref[subsection:II.6.6-fr]{6.6} & Application : critères d'amplitude & 130\\
\hyperref[subsection:II.6.7-fr]{6.7} & Le théorème de Chevalley & 135\\
\hyperref[section:II.7-fr]{\textsection7} & Critères valuatifs & 138\\
\hyperref[subsection:II.7.1-fr]{7.1} & Rappels sur les anneaux de valuation & 138\\
\hyperref[subsection:II.7.2-fr]{7.2} & Critère valuatif de séparation & 141\\
\hyperref[subsection:II.7.3-fr]{7.3} & Critère valuatif de propreté & 143\\
\hyperref[subsection:II.7.4-fr]{7.4} & Courbes algébriques et corps de fonctions de dimension 1 & 148\\
\hyperref[section:II.8-fr]{\textsection8} & Schémas éclatés; cônes projetants; fermeture projective & 152\\
\hyperref[subsection:II.8.1-fr]{8.1} & Préschémas éclatés & 152\\
\hyperref[subsection:II.8.2-fr]{8.2} & Résultats préliminaires sur la localisation dans les anneaux gradués & 157\\
\hyperref[subsection:II.8.3-fr]{8.3} & Cônes projetants & 162\\
\hyperref[subsection:II.8.4-fr]{8.4} & Fermeture projective d'un fibré vectoriel & 168\\
\hyperref[subsection:II.8.5-fr]{8.5} & Comportements fonctoriels & 169\\
\hyperref[subsection:II.8.6-fr]{8.6} & Un isomorphisme canonique pour les cônes épointés & 171\\
\end{longtable}
\oldpage[II]{215}
\markboth{TABLE DES MATIÈRES}{TABLE DES MATIÈRES}
\begin{longtable}{@{}p{0.09\textwidth}p{0.69\textwidth}r@{}}
\hyperref[subsection:II.8.7-fr]{8.7} & Éclatement des cônes projetants & 173\\
\hyperref[subsection:II.8.8-fr]{8.8} & Faisceaux amples et contractions & 177\\
\hyperref[subsection:II.8.9-fr]{8.9} & Le critère d'amplitude de Grauert : énoncé & 182\\
\hyperref[subsection:II.8.10-fr]{8.10} & Le critère d'amplitude de Grauert : démonstration & 184\\
\hyperref[subsection:II.8.11-fr]{8.11} & Unicité des contractions & 189\\
\hyperref[subsection:II.8.12-fr]{8.12} & Faisceaux quasi-cohérents sur les cônes projetants & 191\\
\hyperref[subsection:II.8.13-fr]{8.13} & Fermeture projective de sous-faisceaux et de sous-schémas fermés & 195\\
\hyperref[subsection:II.8.14-fr]{8.14} & Compléments sur les faisceaux associés aux $\mathcal{S}$-Modules gradués & 197\\
\multicolumn{2}{@{}l}{Bibliographie (suite)} & 205\\
\multicolumn{2}{@{}l}{Index des notations} & 207\\
\multicolumn{2}{@{}l}{Index terminologique} & 209\\
\multicolumn{2}{@{}l}{Errata et Addenda (Liste 1)} & 217\\
\end{longtable}

\begin{flushright}
\emph{Reçu le 22 juillet 1960.}
\end{flushright}
\endgroup
